% ===========================================================
%  第 1 章 行列式 —— 高等代数【深化】拓展框（主控撰写）
%  对应 OUTLINE.md §2「深化清单」1.1 – 1.6
%  全部用 \begin{fdeep}{标题} 灰底框；位置由主控嵌入正文相应考点旁
% ===========================================================

% ---------- 1.1 行列式的公理化本质（对应 OUTLINE 1.1）----------
\begin{fdeep}{行列式的本质：反对称多重线性函数}
考研教材里行列式有"六条性质"，一条条背既容易漏，也容易记错方向。
高等代数把这六条压缩成\textbf{三条公理}——凡是满足这三条的映射，有且只有一个，
它就是行列式：
\begin{enumerate}[leftmargin=2.2em,itemsep=0.2em]
\item \textbf{对每一行线性}（多重线性）：固定其余各行时，$|A|$ 是某一行的线性函数；
\item \textbf{换行变号}（反对称性）：交换两行，行列式变号；
\item \textbf{单位阵取 1}（规范性）：$|E|=1$。
\end{enumerate}
【反哺点】：其余性质都能从这三条\textbf{推}出来，不必再背。
\begin{itemize}[leftmargin=2.2em,itemsep=0.2em]
\item 两行相同 $\Rightarrow$ 换这两行后行列式不变又变号 $\Rightarrow 2|A|=0\Rightarrow|A|=0$；
\item 某行全为 $0$ $\Rightarrow$ 由该行线性性得 $|A|=0\cdot(\cdots)=0$；
\item 某行乘 $k$ 加到另一行 $\Rightarrow$ 由线性性拆成"原行列式 $+$ 两行成比例的行列式"$=|A|+0$。
\end{itemize}
\textbf{最容易错的一条}：$|kA|=k^{n}|A|$，而\textbf{不是} $k|A|$。
根因就在"对\textbf{每一行}线性"——$n$ 行各出一个 $k$，所以是 $k^{n}$。
把"$k$ 提出来要数行数"记成习惯，这类送分题就不会再错。
\end{fdeep}

% ---------- 1.2 几何意义（对应 OUTLINE 1.2）----------
\begin{fdeep}{几何意义：$|\det A|$ 是体积，符号是定向}
设 $A=(\alpha_1,\dots,\alpha_n)$，把列向量看成 $n$ 个边向量，则
\fml{|\det A|=\text{这 }n\text{ 个列向量张成的平行体的体积}}
而 $\det A$ 的\textbf{符号}记录这组向量相对标准基的\textbf{定向}。
\fml{\det A=0\iff\text{列向量线性相关}\iff\text{平行体"塌缩"到低维}}
【反哺点】：这一条把三个看似无关的章节打通——
行列式为零、向量组线性相关、齐次方程组有非零解、矩阵不可逆、秩 $<n$，
\textbf{说的是同一件事}。做选择题时无论从哪个角度切入都能互相验证，
也解释了为什么"秩"是贯穿全书的中心概念。
\end{fdeep}

% ---------- 1.3 乘法定理（对应 OUTLINE 1.3）----------
\begin{fdeep}{乘法定理 $|AB|=|A||B|$ 与"一证多得"}
这个定理是以下一串结论的\textbf{共同源头}：
\fml{|AB|=|A||B|,\qquad |A^{k}|=|A|^{k},\qquad |A^{-1}|=|A|^{-1}}
\fml{|A^{*}|=|A|^{n-1},\qquad A\sim B\Rightarrow|A|=|B|}
\textbf{证明思路（分块初等变换）}：构造 $2n$ 阶分块矩阵
\fml{M=\begin{pmatrix}E&A\\ B&O\end{pmatrix}}
第一步，作分块初等变换"第二块行 $\leftarrow$ 第二块行 $-B\times$ 第一块行"，
这不改变行列式的值，于是
\fml{M\longrightarrow\begin{pmatrix}E&A\\ O&-BA\end{pmatrix}}
第二步，按 Laplace 定理沿前 $n$ 行展开：
前 $n$ 行中只有左上角的 $E$ 能取到非零的 $n$ 阶子式，
而把 $E$ 所在的行列调到左上角共需 $n^{2}$ 次对换，故有符号因子 $(-1)^{n^{2}}=(-1)^{n}$，于是
\fml{\begin{vmatrix}E&A\\ O&-BA\end{vmatrix}=(-1)^{n}\,|{-BA}|=(-1)^{n}\cdot(-1)^{n}|AB|=|A||B|}
【反哺点】：这里用到的两件事——"分块初等变换不改变行列式"与"Laplace 展开"——
正好是 9 讲里两个看似孤立的技巧，它们在这里合起来一次解决了乘法定理。
【反哺点】：不必记忆五个孤立公式——它们都是同一个定理的推论。
尤其 $A\sim B\Rightarrow|A|=|B|$ 直接来自 $|P^{-1}AP|=|P^{-1}||A||P|=|A|$，
这正是"相似矩阵行列式相等"的根，也是第 8 章判断相似的必要条件之一。
\end{fdeep}

% ---------- 1.4 按行展开与代数余子式符号（对应 OUTLINE 1.4）----------
\begin{fdeep}{代数余子式的符号 $(-1)^{i+j}$ 从哪来}
把 $n$ 阶行列式按第 $i$ 行展开：
\fml{|A|=\sum_{j=1}^{n}a_{ij}A_{ij},\qquad A_{ij}=(-1)^{i+j}M_{ij}}
那个 $(-1)^{i+j}$ 不是"规定"，而是\textbf{把 $a_{ij}$ 搬到左上角所付出的换行换列代价}：
要把 $a_{ij}$ 移到 $(1,1)$ 位置，需要把它上方的 $i-1$ 行逐行下移、
左方的 $j-1$ 列逐列右移，共交换 $i+j-2$ 次，故符号为 $(-1)^{i+j-2}=(-1)^{i+j}$。
【反哺点】：符号一旦记错，整题归零，而这是最高频的送分变送命点。
理解来源后在考场上可用"$(1,1)$ 位置是正、棋盘式交错"快速自查。
另外记住\textbf{异行异列配错为零}的恒等式：
\fml{\sum_{k=1}^{n}a_{ik}A_{jk}=\begin{cases}|A|,&i=j\\ 0,&i\ne j\end{cases}}
这正是伴随矩阵 $AA^{*}=|A|E$ 的来源，直接服务第 2 章。
\end{fdeep}

% ---------- 1.5 降阶技巧的统一视角（对应 OUTLINE 1.5）----------
\begin{fdeep}{各类降阶技巧为什么都有效}
9 讲给出了加边法、递推法、拆项法、升阶法等套路。它们并非各自独立的技巧，
而是\textbf{多重线性}这同一原理的不同用法：
\begin{itemize}[leftmargin=2.2em,itemsep=0.25em]
\item \textbf{拆项法}——直接使用"对某一行线性"：把一行拆成两行之和，行列式拆成两个行列式之和；
\item \textbf{递推法}——利用行列式结构沿阶数下降，同一线性关系在相邻阶之间重复出现；
\item \textbf{加边（升阶）法}——把 $n$ 阶"升"成 $n+1$ 阶，使新行列式出现可利用的整行/整列公因子，
      再用一次行列式性质把多出来的边"消"掉，本质是\textbf{先用规范性固定尺度}；
\item \textbf{"打洞"消元}——用"某行乘 $k$ 加到另一行不改变行列式"，制造大量零元以启用展开。
\end{itemize}
【反哺点】：遇到没见过的高阶行列式时，不要盲目试套路，先问自己
"这个行列式对哪一行、哪一列具有可利用的线性结构"，
四个方法中总有一个能对上号。这比背题型更能应付新题。
\end{fdeep}

% ---------- 1.6 Vandermonde 与缺项型（对应 OUTLINE 1.6）----------
\begin{fdeep}{Vandermonde 行列式与缺项型变式}
$n$ 阶 Vandermonde 行列式
\fml{V_{n}=\begin{vmatrix}1&1&\cdots&1\\ x_{1}&x_{2}&\cdots&x_{n}\\ \vdots&\vdots&&\vdots\\ x_{1}^{n-1}&x_{2}^{n-1}&\cdots&x_{n}^{n-1}\end{vmatrix}=\prod_{1\le i<j\le n}(x_{j}-x_{i})}
\textbf{缺项型}：若第一行不是全 $1$，而是 $x_{1}^{k},x_{2}^{k},\dots$（缺了某一次幂），
则可利用\textbf{把缺的那一行补上并作差}，或直接用"按列线性性拆项 + 行列式为零项剔掉"的办法处理，
例如含 $x_{i}^{n}$ 与 $x_{i}^{n-2}$ 而缺 $x_{i}^{n-1}$ 的情形，可用
\fml{\begin{vmatrix}1&1&1\\ x_{1}&x_{2}&x_{3}\\ x_{1}^{3}&x_{2}^{3}&x_{3}^{3}\end{vmatrix}=(x_{1}+x_{2}+x_{3})\prod_{i<j}(x_{j}-x_{i})}
这类结论\textbf{不要死记}，掌握"补行相减 / 拆项消零"两条通道即可现场推出。
【反哺点】：考研高频考 Vandermonde 及其变式，考场上现推比背公式更可靠，
且能避免把缺项型的系数记错。
\end{fdeep}
